% Options for packages loaded elsewhere
\PassOptionsToPackage{unicode}{hyperref}
\PassOptionsToPackage{hyphens}{url}
\documentclass[
]{article}
\usepackage{xcolor}
\usepackage[margin=1in]{geometry}
\usepackage{amsmath,amssymb}
\setcounter{secnumdepth}{5}
\usepackage{iftex}
\ifPDFTeX
  \usepackage[T1]{fontenc}
  \usepackage[utf8]{inputenc}
  \usepackage{textcomp} % provide euro and other symbols
\else % if luatex or xetex
  \usepackage{unicode-math} % this also loads fontspec
  \defaultfontfeatures{Scale=MatchLowercase}
  \defaultfontfeatures[\rmfamily]{Ligatures=TeX,Scale=1}
\fi
\usepackage{lmodern}
\usepackage{setspace}
\onehalfspacing % 1.5 line spacing
\ifPDFTeX
  % pdfLaTeX fallback: closest metric-compatible substitute for
  % body font "Times New Roman", serif (from style.css)
  \usepackage{newtx}
\else
  % xetex/luatex font selection: match style.css exactly
  \usepackage{fontspec}
  \setmainfont{Times New Roman} % body { font-family: "Times New Roman", serif; }
  \newfontfamily\headingfont{Georgia} % h1-h6 { font-family: Georgia, serif; }
\fi

% --- Heading font (Georgia, per style.css h1-h6 rule) ---
\usepackage{titlesec}
\ifPDFTeX\else
  \titleformat*{\section}{\headingfont\Large\bfseries}
  \titleformat*{\subsection}{\headingfont\large\bfseries}
  \titleformat*{\subsubsection}{\headingfont\normalsize\bfseries}
\fi

% Use upquote if available, for straight quotes in verbatim environments
\IfFileExists{upquote.sty}{\usepackage{upquote}}{}
\IfFileExists{microtype.sty}{% use microtype if available
  \usepackage[]{microtype}
  \UseMicrotypeSet[protrusion]{basicmath} % disable protrusion for tt fonts
}{}
\makeatletter
\@ifundefined{KOMAClassName}{% if non-KOMA class
  \IfFileExists{parskip.sty}{%
    \usepackage{parskip}
  }{% else
    \setlength{\parindent}{0pt}
    \setlength{\parskip}{6pt plus 2pt minus 1pt}}
}{% if KOMA class
  \KOMAoptions{parskip=half}}
\makeatother
\usepackage{longtable,booktabs,array}
\usepackage{calc} % for calculating minipage widths
% Correct order of tables after \paragraph or \subparagraph
\usepackage{etoolbox}
\makeatletter
\patchcmd\longtable{\par}{\if@noskipsec\mbox{}\fi\par}{}{}
\makeatother
% Allow footnotes in longtable head/foot
\IfFileExists{footnotehyper.sty}{\usepackage{footnotehyper}}{\usepackage{footnote}}
\makesavenoteenv{longtable}
\usepackage{graphicx}
\makeatletter
\newsavebox\pandoc@box
\newcommand*\pandocbounded[1]{% scales image to fit in text height/width
  \sbox\pandoc@box{#1}%
  \Gscale@div\@tempa{\textheight}{\dimexpr\ht\pandoc@box+\dp\pandoc@box\relax}%
  \Gscale@div\@tempb{\linewidth}{\wd\pandoc@box}%
  \ifdim\@tempb\p@<\@tempa\p@\let\@tempa\@tempb\fi% select the smaller of both
  \ifdim\@tempa\p@<\p@\scalebox{\@tempa}{\usebox\pandoc@box}%
  \else\usebox{\pandoc@box}%
  \fi%
}
% Set default figure placement to htbp
\def\fps@figure{htbp}
\makeatother
\setlength{\emergencystretch}{3em} % prevent overfull lines
\providecommand{\tightlist}{%
  \setlength{\itemsep}{0pt}\setlength{\parskip}{0pt}}
\usepackage{booktabs}
\usepackage{caption}
\usepackage{longtable}
\usepackage{colortbl}
\usepackage{array}
\usepackage{anyfontsize}
\usepackage{multirow}
\usepackage{bookmark}
\IfFileExists{xurl.sty}{\usepackage{xurl}}{} % add URL line breaks if available
\urlstyle{same}
\hypersetup{
  pdftitle={NYC Arrest Trends, 2012 - 2024},
  pdfauthor={NYC Data Team},
  hidelinks,
  pdfcreator={LaTeX via pandoc}}

% --- Council logo, top-right header on every page ---
\usepackage{fancyhdr}
\pagestyle{fancy}
\fancyhf{}
\renewcommand{\headrulewidth}{0pt}
\fancyhead[R]{
\includegraphics[width=1.5in]{../assets/data-council-logo.png}}
\setlength{\headheight}{0.6in} % increase if the logo still overlaps the text

\title{NYC Arrest Trends, 2012 - 2024}
\usepackage{etoolbox}
\makeatletter
\providecommand{\subtitle}[1]{% add subtitle to \maketitle
  \apptocmd{\@title}{\par {\large #1 \par}}{}{}
}
\makeatother
\subtitle{Assessing the Relationship between Neighborhood Demographics
and NYPD Arrests in New York City}
\author{NYC Data Team}
\date{August 21, 2026}

\begin{document}
\maketitle

{
\setcounter{tocdepth}{2}
\tableofcontents
}
\section{Introduction and Key
Findings}\label{introduction-and-key-findings}

In 2024, the NYPD made 260,457 arrests across the five boroughs.
Majority of these arrests were of Black and Hispanic New Yorkers,
despite only comprising for 22\% and 27\% of the population,
respectively. Though this citywide disparity in arrest rates is
well-documented, this analysis uncovers trends in arrest incidence at a
more granular level, examining associations between demographic factors
such as race and socioeconomic status and arrest rates. We rely on the
\href{https://data.cityofnewyork.us/Public-Safety/NYPD-Arrests-Data-Historic-/8h9b-rp9u/about_data}{NYPD's
historic arrest data}, available on NYC Open Data. For demographic
indicators, we rely on values derived from the
\href{https://www.census.gov/data/developers/data-sets/acs-5year.html}{American
Community Survey 5-year Data}. To reliably coalesce data from these two
sources, we limit our temporal range from 2012 - 2024, providing summary
statistics with the latest available year in our sample (2024).

The key findings of this analysis can be summarized as follows:

\begin{itemize}
\item
  In the years leading up to the COVID-19 pandemic, the city saw a
  drastic reduction in the number of arrests, with particularly sharp
  declines in predominantly Black and Hispanic CTs relative to
  predominantly White CTs.
\item
  Since the pandemic, arrests have increased, with Black and Hispanic
  CTs experiencing particularly large increases in arrests relative to
  predominantly White CTs.
\item
  Black and Hispanic New Yorkers are overrepresented in city-wide
  arrests, comprising only 22\% and 27\% of NYC residents but accounting
  for 47\% and 36\% of all arrests, respectively. In contrast, White New
  Yorkers, who represent 34\% of the city's population, account for only
  10\% of all arrests.
\item
  Such disparities exist at both a city-wide and local level, with Black
  New Yorkers accounting for a greater proportion of arrests than they
  do as a share of the population in 87\% of the city's census tracts
  (CTs).
\item
  CTs that are predominantly Black or predominantly Hispanic/ Latino
  generally have higher rates of arrests than do predominantly White
  CTs.
\item
  The share of the CT population that is Black is positively correlated
  with the number of arrests in CT. However, this relationship is only
  positive in predominantly White areas. In CTs in which the plurality
  of residents are Black, the number of arrests decreases as the Black
  share of the population increases.
\item
  The proportion of arrests made in predominantly Black CTs that are for
  index crimes as opposed to lower-level offenses is smaller relative to
  the proportion of these arrests made in predominantly White CTs -
  implying that arrests in predominantly Black CTs are more likely to be
  made for lesser offenses, though the difference is quite small.
\end{itemize}

\section{Temporal Trends in NYC
Arrests}\label{temporal-trends-in-nyc-arrests}

\begin{figure}[h]

{\centering \includegraphics[width=0.87\textwidth]{13_pdf_write_up_files/figure-latex/arrests-covid-1} 

}

\caption{Arrests leading up to and following the COVID-19 Pandemic}\label{fig:arrests-covid}
\end{figure}

Prior to the COVID-19 pandemic, arrests rates were declining steadily
throughout the city. Figure \ref{fig:arrests-covid} displays the mean number of arrests per
census tract by predominant ethnicity between 2012 and 2024. Between
2012 and 2024, arrests were rapidly decreasing, with particularly sharp
reductions occurring in predominantly Black and Hispanic CTs. In the
years since 2020, arrests in all census tracts have risen steadily.

In our regression analysis (Table \ref{tab:tab-trend-mods}), we show that leading up to the
COVID-19 pandemic, predominantly Hispanic and Black communities were
experiencing particularly sharp reductions in arrests compared to
predominantly White communities (\emph{Temporal Trend x Pred. Black} and
\emph{Temporal Trend x Pred. Hispanic}, respectively). While all CTs saw
an increase in arrests in the post-pandemic period, the regression
results also imply that these increases were especially pronounced in
predominantly Black and Latino communities (\emph{Post Covid Temporal
Trend x Pred. Black} and \emph{Post-Covid Temporal trend x Hispanic},
respectively). These findings indicate that while Black and Latino
communities experienced the greatest reductions in arrests in the years
leading up to the pandemic, they have also experienced the worst
reversals of those trend in the years since.

\begin{table}[t]
\caption{\label{tab:tab-trend-mods}Estimating Trends in Arrest Rates Before and After the COVID-19 Pandemic} 
\fontsize{12.0pt}{14.0pt}\selectfont
\scalebox{0.85}{%
\begin{tabular*}{\linewidth}{@{\extracolsep{\fill}}lcc}
\toprule
  & (1) & (2) \\ 
\midrule\addlinespace[2.5pt]
Post Covid & \textbf{-0.404} & \textbf{-0.393} \\ 
 & (0.019) & (0.018) \\ 
Temporal Trend (year - 2020) & \textbf{-0.074} & \textbf{-0.074} \\ 
 & (0.003) & (0.003) \\ 
Post Covid Temporal Trend & \textbf{0.213} & \textbf{0.229} \\ 
 & (0.008) & (0.008) \\ 
Post Covid Temporal Trend x Pred. Asian & \textbf{0.043} & \textbf{0.031} \\ 
 & (0.014) & (0.014) \\ 
Post Covid Temporal Trend x Pred. Black & \textbf{0.053} & \textbf{0.051} \\ 
 & (0.011) & (0.011) \\ 
Post Covid Temporal Trend x Pred. Hispanic & \textbf{0.061} & \textbf{0.057} \\ 
 & (0.011) & (0.011) \\ 
Post Covid x Pred. Asian & -0.038 & -0.025 \\ 
 & (0.038) & (0.036) \\ 
Post Covid x Pred. Black & -0.040 & -0.030 \\ 
 & (0.025) & (0.024) \\ 
Post Covid x Pred. Hispanic & -0.014 & -0.002 \\ 
 & (0.025) & (0.024) \\ 
Temporal Trend x Pred. Asian & -0.004 & -0.006 \\ 
 & (0.006) & (0.006) \\ 
Temporal Trend x Pred. Black & \textbf{-0.029} & \textbf{-0.034} \\ 
 & (0.004) & (0.004) \\ 
Temporal Trend x Pred. Hispanic & \textbf{-0.029} & \textbf{-0.032} \\ 
 & (0.004) & (0.004) \\ 
Pred. Asian & -0.014 & -0.013 \\ 
 & (0.033) & (0.032) \\ 
Pred. Black & -0.061 & \textbf{-0.085} \\ 
 & (0.040) & (0.042) \\ 
Pred. Hispanic & \textbf{-0.074} & \textbf{-0.095} \\ 
 & (0.035) & (0.035) \\ 
Population (logged) & \textbf{0.111} & \textbf{0.118} \\ 
 & (0.051) & (0.052) \\ 
Median Income (logged) & \textbf{-0.134} & \textbf{-0.125} \\ 
 & (0.027) & (0.027) \\ 
Wealth Inequality (Gini Index) & -0.205 & -0.190 \\ 
{} & {(0.107)} & {(0.107)} \\ 
Num.Obs. & 27462 & 27462 \\ 
R2 & 0.907 & 0.235 \\ 
R2 Adj. & 0.898 & 0.220 \\ 
R2 Within & 0.427 & 0.161 \\ 
R2 Within Adj. & 0.427 & 0.161 \\ 
AIC & 30467.9 & 247916.5 \\ 
BIC & 49736.9 & 267185.5 \\ 
RMSE & 0.39 & 79.93 \\ 
Std.Errors & by: true\_boroct & by: true\_boroct \\ 
FE: true\_boroct & X & X \\ 
\bottomrule
\end{tabular*}%
}
% \begin{tablenotes}
% \small
% \item \textbf{Bold} indicates $p < 0.05$ (i.e., $|$coefficient / SE$| \geq 1.96$).
% \end{tablenotes}
\end{table}


\clearpage

\section{Geographic and Demographic Trends in NYPD
Arrests}\label{geographic-and-demographic-trends-in-nypd-arrests}

\begin{figure}

{\centering \includegraphics[width=0.85\textwidth]{13_pdf_write_up_files/figure-latex/map-arrests-2024-1} 

}

\caption{Total Arrests by Census Tract, 2024}\label{fig:map-arrests-2024}
\end{figure}

\begin{figure}

{\centering \includegraphics[width=0.85\textwidth]{13_pdf_write_up_files/figure-latex/map-ethnicity-2024-1} 

}

\caption{Predominant ethnicity by census tract, 2024}\label{fig:map-ethnicity-2024}
\end{figure}

Of the more than 260,000 arrests made by the NYPD in 2024, a collective
83\% were of Black or Hispanic New Yorkers. The two maps above display
the total number of arrests in 2024 by census tract (Figure \ref{fig:map-arrests-2024}) and the
predominant ethnicity in each census tract (Figure \ref{fig:map-ethnicity-2024}). Comparing the two,
we can see that aside from Lower Manhattan, arrests are concentrated
within predominantly Black and predominantly Hispanic/ Latino
neighborhoods. Predominantly White census tracts seem to have noticeably
fewer arrests on average. However, we see that some of the census tracts
with the highest concentrations of Black New Yorkers, such as CT 998 in
Canarsie, which is 80\% Black, or CT 1208 in New Lots, which is 93\%
Black, have fewer arrests than average (29 and 25 arrests, respectively,
compared to a citywide average of 114).

Table \ref{tab:overrep} below shows that these disparities exist at both the citywide
and local levels. In 87\% of census tracts, the share of Black New
Yorkers who are arrested is greater than their share of the tract's
population. This is true for those of Hispanic/ Latino origin in 64\% of
tracts and for those of Asian descent in 28\% of tracts. The
corresponding map shows these relationships visually.

\begin{longtable}[]{@{}cc@{}}
\toprule\noalign{}
Race/ Ethnicity & Prop. of CTs Overrepresented in Arrests \\
\midrule\noalign{}
\endhead
\bottomrule\noalign{}
\endlastfoot
Black & 87\% \\
White & 22\% \\
Asian & 28\% \\
Hispanic & 64\% \\
\caption{Proportion of Census Tracts in which each racial/ ethnic
category is overrepresented in arrests relative to their share of the
population.}
\label{tab:overrep}
\end{longtable}

\begin{figure}

{\centering \includegraphics{13_pdf_write_up_files/figure-latex/map-overrep-1} 

}

\caption{By ethnicity, the extent to which each group is overrepresented in CT-level arrests compared to its share of the CT population. The value displayed is the proportion each group makes up of arrests minus the proportion they make up of the CT population. Higher values (blue) indicate that the proportion of those arrested from a given ethnicity is higher than the proportion that ethnicity makes up in the population. Lower values (red) indicate that the reverse.}\label{fig:map-overrep}
\end{figure}

The extent to which CT demographic factors reflect arrest rates is
complex. Across all of the city's census tracts, there is a general
positive and statistically significant association between the share of
the population that is Black and the number of arrests seen in those
communties. Similarly, CTs that are predominantly Black tend to have
significantly more arrests than CTs that are predominantly White.
However, when we disaggregate this analysis and examine the relationship
between the Black population share and arrest rates separately according
to the CT's predominant ethnicity, we observe differential effects. In
predominantly White CTs, there is a positive relationship that mirrors
the citywide averages. Census tracts with greater proportions of Black
residents tend to have higher arrest rates. In contrast, this
relationship becomes negative in predominantly Black communities. There,
a negative association exists between the share of the population that
is Black and the number of arrests that take place in that census tract.
Figure \ref{fig:arrests-black-share} displays these two relationships.

\begin{figure}

{\centering \includegraphics[width=0.8\textwidth]{13_pdf_write_up_files/figure-latex/arrests-black-share-1} 

}

\caption{Black share of CT population and arrest rates, by predominant ethnicity}\label{fig:arrests-black-share}
\end{figure}

\clearpage

\section{Index Crimes}\label{index-crimes}

To gain an understanding of the crimes for which individuals are being
arrested, we examine the number of index crime arrests per census tract.
Index crimes are those which involve serious offenses, such as murder,
robbery, rape, or assault. We track the proportion of arrests that are
made for index crimes as opposed to all other arrests in a census tract.
Figure \ref{fig:map-index-crimes} visualizes the distribution of this proportion around New
York City. Unlike in the previous map, which displayed total arrests,
the spatial distribution of this data is significantly more even, though
several hotspots of index crimes are apparent near Prospect Park and
Northeastern Queens, for example.

\begin{figure}

{\centering \includegraphics[width=0.8\textwidth]{13_pdf_write_up_files/figure-latex/map-index-crimes-1} 

}

\caption{The proportion of 2024 arrests in each census tract that are for index crimes. The value displayed in each CT is the number of arrests for index crimes divided by the total number of arrests in that census tract.}\label{fig:map-index-crimes}
\end{figure}

In Table \ref{tab:tab-index-mod}, we statistically examine whether the relationship
between census tract demographics and index crime rates is similar to
that between demographics and total arrests. Using predominantly White
tracts as a reference category, we see that in predominantly Black
tracts, there are fewer index crime arrests relative to arrests for
lower-level offenses. In predominantly Hispanic and Asian communities,
there are comparatively more arrests made for index crimes relative to
lesser offenses. Though these estimates are statistically significant,
the actual magnitude of these differences is fairly small - just 0.02,
-0.01, and 0.018 for predominantly Asian, Black, and Hispanic
neighborhoods, respectively.

\begin{table}[t]
\caption{\label{tab:tab-index-mod}Estimating the Relationship between Tract Demographics and Index Crime Rates} 
\fontsize{12.0pt}{14.0pt}\selectfont
\scalebox{0.85}{
\begin{tabular*}{\linewidth}{@{\extracolsep{\fill}}lc}
\toprule
  & (1) \\ 
\midrule\addlinespace[2.5pt]
Pred. Asian & \textbf{0.023} \\ 
 & (0.005) \\ 
Pred. Black & \textbf{-0.010} \\ 
 & (0.004) \\ 
Pred. Hispanic & \textbf{0.018} \\ 
 & (0.005) \\ 
Population (logged) & \textbf{0.018} \\ 
 & (0.004) \\ 
Median Income (logged) & -0.002 \\ 
 & (0.004) \\ 
Total Arrests (logged) & 0.001 \\ 
 & (0.002) \\ 
Wealth Inequality (logged) & -0.000 \\ 
{} & {(0.011)} \\ 
Num.Obs. & 27466 \\ 
R2 & 0.106 \\ 
R2 Adj. & 0.106 \\ 
R2 Within & 0.015 \\ 
R2 Within Adj. & 0.014 \\ 
AIC & -33127.6 \\ 
BIC & -32963.2 \\ 
RMSE & 0.13 \\ 
Std.Errors & by: true\_boroct \\ 
FE: yr & X \\ 
\bottomrule
\end{tabular*}
}
% \begin{tablenotes}
% \small
% \item \textbf{Bold} indicates $p < 0.05$ (i.e., $|$coefficient / SE$| \geq 1.96$).
% \end{tablenotes}
\end{table}

\end{document}
